\chapter{A Concrete Naming Certificate for $\ComputableRealUp$}
\label{ch:concrete-instances-computablereal-namecert}

The $\ComputableRealUp$ interface names effective real presentations by finite
BEDC packets rather than by completed real objects. Following the
Bishop--Bridges constructive-analysis treatment of explicit moduli and the
Weihrauch computable-analysis treatment of names as effective approximation
procedures, the packet below records a dyadic approximation program only
through certified finite observation windows, a unary Cauchy modulus, dyadic
endpoint rows, a regular-sequence handoff, and a $\RealUp$ seal boundary.

\begin{definition}[ComputableReal source packet]
\label{def:computable-real-source-packet}
A $\ComputableRealUp$ source packet is a finite $\mathsf{BHist}$ record
$$
X=(s,\mu,d,r,\ell,h,c,p,n)
$$
with the following visible rows:
\begin{enumerate}
\item $s$ is a $\StreamNameUp$ observation schedule naming the certified
finite windows through which the approximation program is read;
\item $\mu$ is a unary $\CauchyModulusUp$ row naming, for each requested
precision row, the finite threshold after which observations are compared;
\item $d$ is the family of $\DyadicRatCoreUp$ packets read on the scheduled
windows selected by $s$ and $\mu$;
\item $r$ is the $\RegSeqRatUp$ regularity handoff ledger built from those
dyadic packets and modulus rows;
\item $\ell$ is the $\RealUp$ seal-source row that consumes the regular
sequence surface without selecting a completed real value;
\item $h$ records the componentwise $\hsame$ rows for stream windows, modulus
rows, dyadic packets, regularity rows, and the seal source;
\item $c$ records the $\Cont$ routes through which certified finite windows
are read by downstream consumers;
\item $p$ is the $\Pkg$ provenance row naming the
$\StreamNameUp$, $\CauchyModulusUp$, $\DyadicRatCoreUp$, $\RegSeqRatUp$, and
$\RealUp$ dependency certificates;
\item $n$ is the local $\NameCert_{\ComputableRealUp}$ registration row.
\end{enumerate}
The carrier predicate is
$\mathsf{ComputableRealSource}(s,\mu,d,r,\ell,h,c,p,n)$. Its classifier
compares two accepted packets by running the existing regular-sequence
classifier on every certified finite window and transporting all displayed
rows componentwise by $\hsame$ and $\Cont$. It does not quotient completed
reals or appeal to host real equality.
\end{definition}

\begin{theorem}[ComputableReal NameCert obligation surface]
\label{thm:computable-real-namecert-obligation-surface}
The local $\NameCert_{\ComputableRealUp}$ obligation surface consists of five
rows:
\begin{enumerate}
\item carrier habitation: the carrier is inhabited by the constant dyadic
stream whose every certified window reads the same $\DyadicRatCoreUp$ zero
packet, with the unary modulus selecting the first threshold row;
\item classifier transport: accepted packets compare only by pointwise
$\hsame$ transport of the displayed $\StreamNameUp$, $\CauchyModulusUp$,
$\DyadicRatCoreUp$, $\RegSeqRatUp$, $\RealUp$, $\Cont$, $\Pkg$, and
$\NameCert$ rows;
\item modulus-window stability: restricting, enlarging, or refining a
certified finite window preserves the late-observation comparisons named by
the same unary modulus row;
\item $\RealUp$ seal handoff: every real-facing read factors through the
$\RegSeqRatUp$ handoff surface and then through the displayed $\RealUp$ seal
source row;
\item ledger coverage: the packet exposes no row beyond the stream schedule,
modulus table, dyadic approximation rows, regularity handoff, seal source,
transport, continuation, provenance, and local naming registration.
\end{enumerate}
These rows are the complete seed obligation for constructive effective real
names over certified dyadic approximations.
\end{theorem}
\begin{proof}
For habitation, take the constant scheduled stream whose dyadic observation
row is the certified zero packet. The unary modulus can select the first
threshold for every precision request, and the regularity handoff reads
identical dyadic endpoints on every finite window.

For classifier transport, project the carrier of
\autoref{def:computable-real-source-packet}. Each visible coordinate is a
finite BEDC row, so the classifier is obtained by applying the
$\RegSeqRatUp$ finite-window classifier after transporting the corresponding
stream, modulus, dyadic, regularity, seal, continuation, provenance, and
local naming rows.

For window stability, compare two certified finite windows by passing to their
displayed common refinement. The unary modulus still names the same threshold
rows, and the dyadic observations beyond those thresholds are the rows already
stored in the packet, so restriction and refinement do not create a new
approximation object.

Finally, the $\RealUp$ read is exactly the seal-source row $\ell$ fed by the
regular-sequence ledger $r$. Since the carrier has no coordinate for a
completed real, a quotient stream, a selected limit, or host real equality,
the listed ledger rows exhaust the local naming surface.
\end{proof}

\begin{theorem}[ComputableReal RealUp handoff]
\label{thm:computable-real-realup-handoff}
Every accepted $\ComputableRealUp$ packet presents a $\RealUp$ seal source
only through its $\RegSeqRatUp$ regularity handoff, dyadic approximation rows,
certified $\StreamNameUp$ windows, unary $\CauchyModulusUp$ threshold rows,
and displayed $\hsame$, $\Cont$, $\Pkg$, and $\NameCert$ provenance.

The handoff exports no completed real carrier, no selected limit, no quotient
of Cauchy names, and no appeal to classical real completeness.
\end{theorem}
\begin{proof}
Project the source packet of
\autoref{def:computable-real-source-packet}. The rows $s$, $\mu$, and $d$
give the scheduled dyadic approximation data, while $r$ is the regular
rational handoff built from those finite observations.

Read the seal-source row $\ell$. It is indexed by the same regularity handoff
and therefore consumes only the finite scheduled windows, dyadic endpoint
rows, modulus thresholds, transport rows, continuation routes, and provenance
already visible in the carrier.

No additional coordinate remains from which a completed real object or
selected limiting value could be read. The $\RealUp$ handoff is therefore a
finite source-surface handoff, not a completeness principle.
\end{proof}

\begin{theorem}[ComputableReal ledger coverage]
\label{thm:computable-real-ledger-coverage}
Every consumer of an accepted $\ComputableRealUp$ packet factors through the
same finite ledger:
$$
(s,\mu,d,r,\ell,h,c,p,n).
$$
The stream schedule, modulus table, dyadic approximation rows, regularity
handoff, seal source, transport rows, continuation routes, package provenance,
and local naming registration cover all readable coordinates of the packet.

Consequently the packet exports no oracle beyond its certified observation
windows, no real-number quotient, no metric completion object, and no hidden
standard bridge.
\end{theorem}
\begin{proof}
Start from the carrier definition. Its tuple lists exactly nine coordinates,
and each coordinate is a finite row over the existing BEDC kernel and
certificate interfaces.

Any stream-facing consumer reads $s$ and the certified windows named by $c$.
Any rate-facing consumer reads $\mu$. Any approximation-facing consumer reads
$d$. Any regular-sequence or real-facing consumer reads $r$ and $\ell$, with
$h$, $p$, and $n$ carrying the transport and provenance required to keep those
reads inside the same packet.

Thus every consumer projection is one of the displayed coordinates or a
projection through the listed transport and continuation rows. Since no other
coordinate exists, external oracles, quotient reals, completion objects, and
standard bridge data are outside the ledger.
\end{proof}

\begin{theorem}[ComputableReal standard bridge boundary]
\label{thm:computable-real-standard-bridge-boundary}
The standard mathematical reading of an accepted $\ComputableRealUp$ packet is
the constructive effective real name consisting of an explicit dyadic
approximation procedure and a unary convergence modulus. The BEDC bridge
boundary is exactly the finite source surface of
\autoref{def:computable-real-source-packet} together with the obligation rows
of \autoref{thm:computable-real-namecert-obligation-surface},
\autoref{thm:computable-real-realup-handoff}, and
\autoref{thm:computable-real-ledger-coverage}.

This bridge boundary is compatible with Bishop--Bridges constructive reals and
Weihrauch-style computable names, but it does not assert classical real
completeness, a quotient of names, extensional equality of completed reals, or
a total comparison operation.
\end{theorem}
\begin{proof}
Read the accepted packet as the finite program interface encoded by
$s$, $\mu$, and $d$: certified windows select dyadic observations, and the
unary modulus selects the finite thresholds at which precision requests are
checked.

Apply \autoref{thm:computable-real-namecert-obligation-surface} to identify
the carrier, pointwise classifier, modulus-window stability, seal handoff, and
ledger coverage rows. Apply
\autoref{thm:computable-real-realup-handoff} to identify the only real-facing
read, and apply \autoref{thm:computable-real-ledger-coverage} to rule out
unlisted consumer coordinates.

The resulting bridge boundary is therefore the standard effective real-name
surface interpreted through finite BEDC rows. The external precedents explain
the intended mathematical reading, while the BEDC packet itself retains only
the displayed approximation, modulus, regularity, seal, transport, and
provenance rows.
\end{proof}

\closureat{ComputableRealUp}{seedStr}

\begin{closurestatus}{\ComputableRealUp}
  \constructivestory{A computable real source is generated by certified finite stream-observation windows, a unary Cauchy modulus, dyadic approximation packets, a regular-sequence handoff, a RealUp seal source, and visible transport and provenance rows.}
  \theoryclosure{\seedClosure}
  \scopeclosed{\autoref{def:computable-real-source-packet}, \autoref{thm:computable-real-namecert-obligation-surface}, \autoref{thm:computable-real-realup-handoff}, and \autoref{thm:computable-real-ledger-coverage} bind the seed carrier, obligation rows, real-seal handoff, and finite ledger coverage for constructive effective real names.}
  \formalstatus{\unformalizedV}
  \bridgestatus{none}
  \origin{human}
  \notclaimed{This chapter does not close classical real completeness, quotient equality of names, total comparison of completed reals, selected limits, or a machine-checked standard bridge.}
  \upgradepath{Formal carrier, classifier, modulus-window stability, RealUp handoff, ledger coverage, and bridge-boundary targets are required before the surface can move beyond seed status.}
\end{closurestatus}
